\vspace{0.5em}\noindent\textbf{{[}FCAJ2026{]} AWS S3 Bucket: The Cloud-Native Storage
``Missing Piece'' for Digital Healthcare Systems}

\vspace{0.5em}\noindent\textbf{Introduction}

During the development of our Digital Healthcare System (an AI-assisted
medical consultation and appointment management platform), after solving
the problems of anti-overlap appointment booking, integrating AI models,
and building the user authentication flow, I faced another major
headache: \textbf{File Data Management and Storage (File Storage)}.

Our system continuously generates a wide variety of file data: *
\textbf{Patients:} Avatars, medical summary reports/prescriptions in PDF
format. * \textbf{Doctors \& Staff:} Professional degrees, medical
practice certificates. * \textbf{AI Models:} Training weight files (.pt,
.bin) up to hundreds of MBs in size, along with consultation chat logs.

Initially, my most convenient habit when working locally was to dump
certificate and avatar files directly into the \texttt{uploads/} folder
of the NestJS Backend server, and commit the AI weight files straight
into Git. However, as we prepared to push the system to Production, I
realized this was a ``disaster'' in terms of security and scalability:
the Git repo became too heavy, the Backend server was burdened with file
storage (Stateful), and the risk of leaking personal medical files was
extremely high.

That was when I shifted to exploring and applying \textbf{Amazon S3
(Simple Storage Service)}.

\vspace{0.5em}\noindent\textbf{Practical Application in the Project}

Instead of turning the Backend server into a storage ``dump,'' I
restructured the entire system to push all file-based resources to an S3
Bucket. Specifically, the strategic application areas included:

\begin{itemize}
\tightlist
\item
  \textbf{Securing Doctor Practice Certificates:} Medical degrees are
  sensitive personal data that absolutely cannot be made public. I
  created a Private S3 Bucket to hold these files. When an Admin needs
  to review a doctor's profile, the NestJS Backend generates an S3
  Presigned URL with a short expiration time. Once expired, the URL is
  automatically invalidated, ensuring absolute security.
\item
  \textbf{AI Model Weights Management:} Weight files (\texttt{model.pt})
  were completely separated from the source code. All weights and
  training checkpoints are stored on S3. When the AI service container
  starts up, a Python script automatically pulls the latest weights from
  S3 into memory to load the model. Version management for the AI model
  became extremely lightweight.
\item
  \textbf{Exporting Medical Reports \& Storing Chat Logs:} Whenever a
  patient completes a consultation or examination session, the system
  automatically exports a PDF report and records the chat log from the
  AI Service. These files are pushed directly to an S3 Bucket for
  downloading or long-term storage.
\end{itemize}

\vspace{0.5em}\noindent\textbf{Outstanding Values When Switching to Amazon S3}

After successfully integrating S3 Buckets into the system, here are the
biggest values the project gained:

\begin{itemize}
\tightlist
\item
  \textbf{Freeing the Backend to a ``Stateless'' state:} The NestJS
  Backend server now purely handles business logic and APIs. Removing
  all static files makes the server much lighter and easy to scale-out
  to multiple instances without worrying about out-of-sync stored files.
\item
  \textbf{Absolute Security for Sensitive Medical Data:} There is no
  longer a risk of exposing direct file paths. Access control through
  IAM Roles and Presigned URLs ensures that only the right people at the
  right time can access medical documents.
\item
  \textbf{Ultimate Reliability \& Optimized Costs:} S3 guarantees data
  durability of up to 99.999999999\% (11 nines). Besides, AWS provides
  5GB of free storage for the first 12 months. Combined with Lifecycle
  Rules (automatically deleting temporary logs or moving old files to
  the low-cost Glacier storage), the cost of maintaining the storage
  infrastructure for the project is practically zero.
\end{itemize}

\vspace{0.5em}\noindent\textbf{Conclusion}

Applying Cloud-Native services like Amazon S3 to manage files is the
core mindset that helps standardize modern system architecture. Instead
of turning the Backend server into a bulky ``warehouse'' and facing a
series of security risks, shift the burden of storage and resource
management to AWS S3.

This not only makes your project more professional and flexible but also
allows you to focus entirely on developing core business features.

\vspace{0.5em}\noindent\textbf{References}

\begin{itemize}
\tightlist
\item
  AWS. ``Amazon Simple Storage Service (S3) User Guide''.
  (\href{https://docs.aws.amazon.com/AmazonS3/latest/userguide/}{Link})
\item
  AWS. ``Amazon S3 Product Overview''.
  (\href{https://aws.amazon.com/s3/}{Link})
\end{itemize}
